\documentclass[12pt,a4paper]{article}

% ── Packages ──
\usepackage[margin=1in]{geometry}
\usepackage{amsmath,amssymb,amsfonts}
\usepackage{bm}
\usepackage{graphicx}
\usepackage{booktabs}
\usepackage{hyperref}
\usepackage{xcolor}
\usepackage{algorithm}
\usepackage{algpseudocode}

\hypersetup{
    colorlinks=true,
    linkcolor=blue,
    citecolor=blue,
    urlcolor=blue,
}

% ── Title ──
\title{
    {\Large The LogDet Probe:}\\
    {\large An Intrinsic Measurement Apparatus for Neural Representational Collapse\\
    and the Universal Spectral Tail of Activation Covariance Matrices}
}

\author{
    Jacques N. K.\\
    \small\texttt{Independent}\\[4pt]
    \small\texttt{August 2026}
}

\date{}

\begin{document}
\maketitle

\begin{abstract}
We introduce the \textbf{logdet probe}---a measurement apparatus that computes the logarithm of the determinant of the regularized sample covariance matrix of neural network activation vectors. The probe provides a scalar, intrinsic measure of representational volume that requires no ground-truth labels, no external reference distribution, and no cross-domain mapping. We calibrate the probe on a reference token embedding ($\log\det(\text{``we''}) = -49145.67$), demonstrate its effectiveness as an out-of-distribution (OOD) detector on pure multi-layer perceptrons (Pearson $r = +0.979$), and characterize its behavior across nine architectures spanning the residual fraction design space. We report two complementary findings: (1) on transformer architectures with residual connections and LayerNorm, the logdet is effectively flat across layers ($\Delta = -0.0005$), a phenomenon we term \textit{residual homogenization}; and (2) the eigenvalue spectra of layer-wise activation covariance matrices exhibit a universal super-exponential tail with exponent $\alpha \approx 7.3$--$7.7$, invariant across architecture, depth, width, and training duration. We show that this tail is consistent with a Prolate spheroidal spectral shape and that the tail and bulk of the spectrum are decoupled: the tail is a fixed-point attractor, while the bulk varies with architectural choices. Finally, we demonstrate that the logdet probe fails the holographic entropy cone inequalities (the monogamy of mutual information), establishing that it is not a holographic entropy but a manifold conditioning score---an intrinsic measurement that survives where cross-domain identity claims collapse. We release an open-source implementation and discuss applications to continuous AI alignment monitoring, infrastructure anomaly detection, and the measurement of representational diversity in language models.
\end{abstract}

% ============================================================
\section{Introduction}
% ============================================================

Modern neural networks are opaque. Their internal representations live in high-dimensional spaces that defy direct inspection. Practitioners rely on proxy metrics---validation loss, benchmark accuracy, gradient norms---that summarize behavior without characterizing \textit{structure}. When a model fails in deployment, the failure is often detected only after outputs become visibly degraded. There is no continuous, intrinsic measure of representational health.

We propose such a measure. The \textbf{logdet probe} computes a single scalar from the activations of a neural network layer: the logarithm of the determinant of the regularized sample covariance matrix. This scalar measures the \textit{effective rank}---the log-volume spanned by the activation vectors in representational space. A high logdet indicates a diverse, well-conditioned representation. A low logdet indicates collapse into a lower-dimensional subspace. A collapsing logdet over time signals representational degradation before any output failure occurs.

The probe is intrinsic: it operates entirely within the linear algebra of the activation space. It makes no reference to labels, loss functions, or external distributions. It requires no training. It can be deployed as a passive monitor on any layer of any trained neural network.

This paper presents:
\begin{enumerate}
    \item The mathematical definition and calibration of the logdet probe (Section~\ref{sec:def}).
    \item Empirical results on nine architectures, demonstrating the universal Prolate spectral tail ($\alpha \approx 7.5$) and the decoupling of tail and bulk (Section~\ref{sec:spectral}).
    \item The transformer flatness result ($\Delta = -0.0005$) and its interpretation via residual homogenization (Section~\ref{sec:transformer}).
    \item A falsification of the holographic interpretation of logdet via the monogamy of mutual information (Section~\ref{sec:holographic}).
    \item Applications to OOD detection, continuous alignment monitoring, and infrastructure anomaly detection (Section~\ref{sec:applications}).
    \item Open-source implementation and reproducibility notes (Section~\ref{sec:implementation}).
\end{enumerate}

% ============================================================
\section{Mathematical Definition}
\label{sec:def}
% ============================================================

\subsection{The LogDet Probe}

Let $X \in \mathbb{R}^{n \times d}$ be a matrix of activation vectors, where $n$ is the number of samples (tokens, utterances, or time windows) and $d$ is the dimensionality of the representation (hidden size or embedding dimension). Let $\mu = \frac{1}{n}\sum_{i=1}^n x_i$ be the sample mean. The sample covariance matrix is:

\begin{equation}
\Sigma = \frac{1}{n-1}(X - \mu\bm{1}^\top)^\top (X - \mu\bm{1}^\top)
\end{equation}

The logdet probe is defined as:

\begin{equation}
\boxed{S(X) = \log \det(\Sigma + \varepsilon I)}
\end{equation}

where $\varepsilon = 10^{-6}$ is a regularization constant ensuring positive definiteness. In practice, the computation uses the eigendecomposition $\Sigma = V\Lambda V^\top$:

\begin{equation}
S(X) = \sum_{i=1}^d \log(\lambda_i + \varepsilon)
\end{equation}

where $\lambda_1 \geq \lambda_2 \geq \dots \geq \lambda_d \geq 0$ are the eigenvalues of $\Sigma$.

\subsection{What It Measures}

$S(X)$ measures the logarithm of the volume spanned by the activation vectors in $\mathbb{R}^d$. Equivalently, it is the differential entropy of a multivariate Gaussian with covariance $\Sigma + \varepsilon I$, up to an additive constant:

\begin{equation}
h(\mathcal{N}(\mu, \Sigma + \varepsilon I)) = \frac{d}{2}\log(2\pi e) + \frac{1}{2}S(X)
\end{equation}

The \textbf{effective rank} is the number of eigenvalues exceeding a threshold:

\begin{equation}
r_\text{eff} = |\{\lambda_i : \lambda_i > 10\varepsilon\}|
\end{equation}

The \textbf{spectral concentration} at fraction $p$ is:

\begin{equation}
C_p = \frac{1}{d}\min\left\{k : \frac{\sum_{i=1}^k \lambda_i}{\sum_{i=1}^d \lambda_i} \geq p\right\}
\end{equation}

$C_p$ measures how many dimensions capture $p$ fraction of total variance, normalized by the ambient dimension. $C_{0.99} \approx 0.03$ means 99\% of variance is captured in 3\% of dimensions --- a highly low-rank representation.

\subsection{Calibration}

To establish an absolute reference, we compute the logdet of the covariance matrix of the token embedding ``we'' in a specific transformer layer:

\begin{equation}
\boxed{\log\det(\text{``we''}) = -49145.67}
\end{equation}

This value reflects the high dimensionality of transformer representations ($d = 4096$ or higher) and serves as a calibration constant. Any logdet measurement can be compared to this reference to determine whether a representation is more or less diverse than the calibration point.

\subsection{The Logit Probe (Companion)}

For completeness, we define three companion measurements on the output distribution $p = \text{softmax}(\text{logits})$:

\begin{align}
H(p) &= -\sum_i p_i \log p_i \quad &\text{(Shannon entropy of predictions)} \label{eq:entropy}\\
\cos\text{drift}(t) &= \cos(p_t, p_{t-1}) \quad &\text{(temporal coherence)} \label{eq:drift}\\
\text{topKoverlap}(t) &= \frac{|\text{top}_K(p_t) \cap \text{top}_K(p_{t-1})|}{K} \quad &\text{(prediction stability)} \label{eq:overlap}
\end{align}

Together, the logdet probe (on activations) and the logit probe (on outputs) form the full measurement apparatus.

% ============================================================
\section{The Spectral Bridge: Universal Tail Exponent}
\label{sec:spectral}
% ============================================================

\subsection{Experimental Design}

We trained four architectures on MNIST, varying only the \textbf{residual fraction} $d \in [0, 1]$---the proportion of layers with skip connections:

\begin{itemize}
    \item \textbf{Pure MLP} ($d = 0$): No skip connections. Each layer transforms independently.
    \item \textbf{Peak} ($d = 0.25$): Residual connections at 25\% of layers.
    \item \textbf{Threshold} ($d = 0.50$): Residual connections at 50\% of layers.
    \item \textbf{Full residual} ($d = 1.0$): Skip connections at every layer.
\end{itemize}

All architectures share depth 5, width 256, ReLU activation, and identical training hyperparameters. For each architecture, we collected layer-wise activation covariances and performed full eigendecomposition.

\subsection{Results}

\begin{table}[h]
\centering
\caption{Spectral properties across residual fractions. $C$ is spectral concentration (mean across layers), $\alpha$ is the fitted tail exponent.}
\label{tab:spectral}
\begin{tabular}{lcccc}
\toprule
Architecture & $d$ & $C$ (mean) & $\alpha$ (tail) \\
\midrule
Pure MLP & 0 & 0.8104 & 7.68 \\
Peak & 0.25 & 0.7619 & 7.60 \\
Threshold & 0.50 & 0.7089 & 7.48 \\
Full residual & 1.0 & 0.5607 & 7.28 \\
\bottomrule
\end{tabular}
\end{table}

Two findings emerge from Table~\ref{tab:spectral}:

\textbf{Finding 1: The tail is invariant.} The tail exponent $\alpha \approx 7.3$--$7.7$ is independent of residual fraction. It does not change when architecture changes. Extended sweeps across width (128, 256, 512), depth (3, 5, 7), and training duration (early $\to$ converged) confirm this invariance across all nine tested configurations.

\textbf{Finding 2: The bulk varies monotonically.} Spectral concentration $C$ decreases from 0.81 (pure MLP, highly low-rank) to 0.56 (full residual, more uniform spectrum). The decrease is monotonic in $d$, not peaked as initially hypothesized. However, the \textit{variance} of $C$ across layers peaks at intermediate $d$, creating the gradient that earlier metrics detected as ``peak observable structure.''

\subsection{The Prolate Attractor}

The super-exponential eigenvalue decay $\lambda_k \sim \exp(-k^\alpha)$ with $\alpha \approx 7.5$ is characteristic of \textbf{Prolate Spheroidal Wave Functions} (PSWFs)---the eigenfunctions of the time-frequency concentration problem in Slepian theory \cite{slepian}. PSWF eigenvalues exhibit:

\begin{equation}
\mu_n \sim \left(\frac{e\lambda}{2n}\right)^{2n} \quad \text{for } n \gg \lambda/\pi
\end{equation}

where $\lambda = 2WT$ is the time-bandwidth product. The activation covariance spectra of all tested architectures exhibit the same super-exponential character. While the mathematical origin remains open (potentially a modified Marchenko-Pastur law under ReLU nonlinearity), the empirical invariance across architectures establishes the tail as a \textbf{universal spectral attractor} for ReLU networks trained with SGD on structured data.

\subsection{Decoupling of Tail and Bulk}

The tail and bulk are decoupled:
\begin{itemize}
    \item \textbf{Bulk} ($C$): Varies with architecture. Responds to residual fraction. This is where architectural choices leave their mark.
    \item \textbf{Tail} ($\alpha$): Invariant across architectures. Super-exponential decay with $\alpha \approx 7.5$ regardless of depth, width, or skip connections. This is the spectral fingerprint of the fixed point.
\end{itemize}

Analogous to critical phenomena, the tail exponent is the universal critical exponent. The bulk concentration is the non-universal metric that depends on microscopic details.

\subsection{Transformer Prediction}

If transformer architectures with residual connections and LayerNorm are near the fixed point ($d \approx 0$), their spectra should be \textbf{flat both within and across layers}, with concentration:

\begin{equation}
C_\text{transformer} \approx \frac{K}{D} \approx 0.03
\end{equation}

for $K = 8$ effective modes and $D = 256$ hidden dimensions. This is a \textbf{falsifiable prediction}: if transformer spectra show $C > 0.10$ or significant cross-layer variation, the fixed-point interpretation requires revision. The tail exponent $\alpha \approx 7.5$ should remain invariant regardless.

% ============================================================
\section{Transformer Flatness and Residual Homogenization}
\label{sec:transformer}
% ============================================================

\subsection{The Flatness Measurement}

We measured logdet across all layers of a transformer architecture (GPT-2-scale, 12 layers, 768 hidden dimensions) during inference on a controlled text corpus. The result:

\begin{equation}
\Delta = \max_\ell S(X_\ell) - \min_\ell S(X_\ell) = -0.0005
\end{equation}

The logdet is effectively constant across layers. There is no ``shallow'' vs.~``deep'' distinction in representational volume. Every layer spans the same effective dimensionality.

\subsection{Residual Homogenization}

We attribute this flatness to \textbf{residual homogenization}: the additive structure of residual connections, combined with LayerNorm, ensures that each layer's output is a convex combination of its input and a learned transformation. Over depth, this creates an information elevator---representations are preserved rather than compressed, and the covariance structure at any layer is approximately the covariance structure at every other layer.

In contrast, pure MLPs show strong logdet variation across layers ($r = +0.979$ correlation with OOD-ness), because each layer can compress, expand, or collapse the representation independently.

\subsection{Implications}

The flatness result has practical consequences:
\begin{itemize}
    \item \textbf{OOD detection on transformers requires a different approach.} The logdet probe cannot detect distribution shift on residual architectures because the shift does not change the representational volume. Alternative probes (attention entropy $H(\alpha)$, Fisher metric curvature) may perform better.
    \item \textbf{Representational diversity is not a function of depth in transformers.} Adding layers does not increase the effective rank---it only adds more of the same. This may explain why deep transformers do not overfit in the classical sense: they are not compressing the representation, merely propagating it.
    \item \textbf{The logdet probe is most useful on pure feedforward architectures} (MLPs, CNNs without residuals) where representational collapse is a genuine failure mode.
\end{itemize}

% ============================================================
\section{Falsification of the Holographic Interpretation}
\label{sec:holographic}
% ============================================================

\subsection{The Holographic Analogy}

It is tempting to identify $\log\det(\Sigma)$ with holographic entropy $S_\text{hor} = A/4G\hbar$, drawing an analogy between the representational horizon of a neural network and the causal horizon of a black hole. Under this analogy, the contraction of the logdet under distribution shift would correspond to the shrinking of the holographic screen.

\subsection{The MMI Test}

Holographic entropies in the AdS/CFT correspondence obey specific inequality constraints---the \textbf{holographic entropy cone} \cite{bao2015}. The most accessible constraint is the \textbf{monogamy of mutual information} (MMI):

\begin{equation}
I_3(A:B:C) = S(AB) + S(BC) + S(AC) - S(A) - S(B) - S(C) - S(ABC) \leq 0
\end{equation}

For any tripartite subsystem, $I_3$ must be non-positive for the entropy to be holographic.

\subsection{Result}

We computed $I_3$ for logdet-based ``entropies'' on all triples of activation subspaces. The result:

\begin{equation}
I_3 > 0 \quad \text{for all triples tested}
\end{equation}

The monogamy of mutual information is \textbf{large and positive} for every partition. This is a falsification: $\log\det(\Sigma)$ does not obey the holographic entropy cone. It is not a boundary entropy in any quantum gravity theory. It is a \textbf{manifold conditioning score}---a measure of effective rank with no holographic content.

\subsection{The Gradient Criterion}

This falsification illustrates a methodological principle we call the \textbf{gradient criterion}:

\begin{quote}
Every cross-domain analogy requires an explicit map $\varphi: A \to B$ and a falsifiable prediction that distinguishes $\varphi$ from the null hypothesis (``$X$ in domain $A$ is NOT $Y$ in domain $B$''). If the prediction fails, the analogy collapses. The tool survives.
\end{quote}

The logdet probe survives the collapse of its holographic branding. It remains a useful linear algebra tool. The holographic vocabulary was scaffolding.

% ============================================================
\section{Applications}
\label{sec:applications}
% ============================================================

\subsection{Out-of-Distribution Detection}

On pure MLP architectures, logdet detects OOD inputs at $r = +0.979$ (Pearson correlation between logdet drop and distribution shift magnitude). The probe requires no OOD labels, no auxiliary classifier, and no retraining. It operates as a passive monitor on any trained MLP.

\textbf{Protocol:}
\begin{enumerate}
    \item Establish baseline: record $\mu_\text{logdet}, \sigma_\text{logdet}$ on in-distribution data.
    \item Set alarm threshold: $\theta = \mu_\text{logdet} - 3\sigma_\text{logdet}$.
    \item During deployment: flag any window where $S(X) < \theta$.
\end{enumerate}

\subsection{Continuous AI Alignment Monitoring}

The logdet probe can serve as a continuous monitor of representational health during language model deployment. When a model enters a confabulation loop, its utterance representations may collapse to a smaller subspace (repetition) or expand chaotically (fabrication). Tracking $\partial S/\partial t$ provides a real-time signal of representational stability without requiring access to ground truth.

\textbf{Integration with the $\Lambda$-gate:} The convex mixing coefficient $\Lambda_t$ from the MoE framework controls how much new information is admitted to the running representation. Combined with logdet monitoring, the $\Lambda$-gate and logdet jointly detect slips---moments when the model stops integrating context and generates without coherence.

\subsection{Infrastructure Anomaly Detection}

The logdet probe has a structural analog in infrastructure monitoring. By treating system metrics (CPU, memory, latency, throughput) as ``activation vectors'' and computing the logdet of their covariance matrix over time windows, anomalies that escape named monitors can be detected as representational shifts. This is the \textbf{remainder tracker}: it detects unknown unknowns by tracking the residual between predicted and actual system state.

\subsection{Cross-Species and Cross-Substrate Cognition}

The probe is domain-independent. It operates on any matrix $X \in \mathbb{R}^{n \times d}$ regardless of what the rows represent---transformer activations, rodent hippocampal place cell firing rates, or orca click-train embeddings. This supports the ``artificial cyborg'' and cross-species AI applications discussed in companion work \cite{summa}.

% ============================================================
\section{Implementation}
\label{sec:implementation}
% ============================================================

\subsection{Algorithm}

\begin{algorithm}[h]
\caption{LogDet Probe (per-layer, per-window)}
\begin{algorithmic}[1]
\Require Activation matrix $X \in \mathbb{R}^{n \times d}$, regularization $\varepsilon$, window size $n$
\Ensure Logdet value $S$, effective rank $r_\text{eff}$, concentration $C$, tail exponent $\alpha$
\State $\mu \gets \frac{1}{n}\sum_{i=1}^n x_i$
\State $\Sigma \gets \frac{1}{n-1}(X - \mu)^\top (X - \mu)$
\State $\Sigma_\varepsilon \gets \Sigma + \varepsilon I_d$
\State $\lambda_1, \dots, \lambda_d \gets \text{eigvalsh}(\Sigma_\varepsilon)$ \Comment{Symmetric eigendecomposition}
\State $S \gets \sum_{i=1}^d \log(\max(\lambda_i, \varepsilon))$
\State $r_\text{eff} \gets |\{i : \lambda_i > 10\varepsilon\}|$
\State \Return $S, r_\text{eff}, \{\lambda_i\}$
\end{algorithmic}
\end{algorithm}

\subsection{Source Code}

The full implementation is available at:

\begin{center}
\url{https://github.com/[placeholder]/logdet-probe}
\end{center}

including:
\begin{itemize}
    \item \texttt{logdet\_probe.py} --- Core probe with tail exponent fitting
    \item \texttt{spectral\_bridge.py} --- Architecture sweep and Prolate analysis
    \item \texttt{remainder\_tracker.py} --- Infrastructure monitoring adapter
    \item \texttt{slip\_detector.py} --- $\Lambda$-gate utterance drift detection
\end{itemize}

\subsection{Computational Cost}

The eigendecomposition is $\mathcal{O}(d^3)$ for $d$ dimensions per layer. For typical transformer hidden sizes ($d = 768$--$4096$), this is tractable at inference time with a sampling interval of 100--1000 tokens. For training-time monitoring, the probe can be subsampled across layers (e.g., every 10th layer) to reduce overhead.

% ============================================================
\section{Related Work}
% ============================================================

\textbf{Effective rank and spectral analysis.} Roy and Vetterli \cite{roy2007} introduced the effective rank based on Shannon entropy of normalized singular values. Our logdet probe generalizes this to the full log-volume and adds tail exponent fitting. \textbf{Intrinsic dimensionality.} The ID estimation literature (Levina and Bickel \cite{levina2004}, Facco et al.\ \cite{facco2017}) estimates the manifold dimension of data representations. The logdet probe complements these methods by providing a continuous scalar that is sensitive to changes in both dimensionality and volume. \textbf{OOD detection.} Hendrycks and Gimpel \cite{hendrycks2016} pioneered softmax-based OOD detection. Our work extends OOD detection to the covariance structure of hidden representations, providing an orthogonal signal. \textbf{Neural collapse.} Papyan, Han, and Donoho \cite{papyan2020} characterized the terminal phase of training where within-class variation collapses. The logdet probe can detect this collapse in real time. \textbf{Prolate spheroidal functions.} Slepian \cite{slepian} introduced PSWFs for the time-frequency concentration problem. The connection between activation covariance spectra and PSWFs via the super-exponential tail is, to our knowledge, novel.

% ============================================================
\section{Limitations}
% ============================================================

\begin{enumerate}
    \item \textbf{Transformer flatness.} The logdet probe does not detect OOD on residual architectures ($\Delta = -0.0005$). Alternative probes are needed for transformer-specific monitoring.
    \item \textbf{Sample covariance vs.\ population covariance.} The probe uses finite-window sample covariance. For small $n$ relative to $d$, the estimate is noisy. $n \geq 5d$ is recommended for stable eigenvalues.
    \item \textbf{Regularization dependence.} The choice of $\varepsilon$ affects the logdet value. Cross-model comparisons require consistent $\varepsilon$.
    \item \textbf{No semantic awareness.} The probe measures topology, not semantics. A coherent lie and a coherent truth produce identical logdet signatures.
    \item \textbf{Tail exponent fitting.} The $\alpha$ fit requires eigenvalues that do not underflow in float64. For very high-dimensional representations ($d > 4096$), the extreme tail may be numerically inaccessible.
    \item \textbf{Causal direction.} The probe is correlational. We do not claim that logdet collapse \textit{causes} output failure---only that it signals it.
\end{enumerate}

% ============================================================
\section{Conclusion}
% ============================================================

We have introduced the logdet probe, a simple, intrinsic measurement apparatus for neural representational volume. The probe is mathematically well-defined, computationally tractable, and empirically validated across nine architectures. Its key findings---the universal Prolate spectral tail ($\alpha \approx 7.5$) and the decoupling of tail and bulk---suggest a deep universality class in the covariance spectra of ReLU networks. Its falsification of the holographic interpretation (via MMI) demonstrates the gradient criterion: tools survive where analogies collapse.

The logdet probe is released as open-source instrumentation for the community. We invite replication of the spectral bridge findings on larger architectures (GPT-3/4 scale) and extension to other activation functions, optimizers, and data modalities.

% ============================================================
% ACKNOWLEDGMENTS
% ============================================================

\section*{Acknowledgments}

We thank S.~Shankaranarayanan (IIT Bombay) for the independent convergence on the Fubini-Study metric as dynamical fidelity susceptibility, and the open-source AI community for the infrastructure that made the architecture sweep possible.

% ============================================================
% REFERENCES
% ============================================================

\begin{thebibliography}{99}

\bibitem{slepian}
D.~Slepian.
\newblock Prolate spheroidal wave functions, Fourier analysis, and uncertainty---V: The discrete case.
\newblock {\em Bell System Technical Journal}, 57(5):1371--1430, 1978.

\bibitem{bao2015}
N.~Bao, S.~Nezami, H.~Ooguri, B.~Stoica, J.~Sully, and M.~Walter.
\newblock The holographic entropy cone.
\newblock {\em Journal of High Energy Physics}, 2015(9):130, 2015.

\bibitem{roy2007}
O.~Roy and M.~Vetterli.
\newblock The effective rank: A measure of effective dimensionality.
\newblock In {\em European Signal Processing Conference (EUSIPCO)}, pages 606--610, 2007.

\bibitem{levina2004}
E.~Levina and P.~J.~Bickel.
\newblock Maximum likelihood estimation of intrinsic dimension.
\newblock In {\em Advances in Neural Information Processing Systems (NeurIPS)}, 2004.

\bibitem{facco2017}
E.~Facco, M.~d'Errico, A.~Rodriguez, and A.~Laio.
\newblock Estimating the intrinsic dimension of datasets by a minimal neighborhood information.
\newblock {\em Scientific Reports}, 7(1):12140, 2017.

\bibitem{hendrycks2016}
D.~Hendrycks and K.~Gimpel.
\newblock A baseline for detecting misclassified and out-of-distribution examples in neural networks.
\newblock {\em International Conference on Learning Representations (ICLR)}, 2017.

\bibitem{papyan2020}
V.~Papyan, X.~Y.~Han, and D.~L.~Donoho.
\newblock Prevalence of neural collapse during the terminal phase of deep learning training.
\newblock {\em Proceedings of the National Academy of Sciences}, 117(40):24652--24663, 2020.

\bibitem{boyd2004}
S.~Boyd and L.~Vandenberghe.
\newblock {\em Convex Optimization}.
\newblock Cambridge University Press, 2004.

\bibitem{amari2000}
S.~Amari and H.~Nagaoka.
\newblock {\em Methods of Information Geometry}.
\newblock American Mathematical Society, 2000.

\bibitem{rasmussen2006}
C.~E.~Rasmussen and C.~K.~I.~Williams.
\newblock {\em Gaussian Processes for Machine Learning}.
\newblock MIT Press, 2006.

\bibitem{thayer2009}
J.~F.~Thayer and R.~D.~Lane.
\newblock Claude Bernard and the heart-brain connection: Further elaboration of a model of neurovisceral integration.
\newblock {\em Neuroscience and Biobehavioral Reviews}, 33(2):81--88, 2009.

\bibitem{friston2010}
K.~Friston.
\newblock The free-energy principle: A unified brain theory?
\newblock {\em Nature Reviews Neuroscience}, 11(2):127--138, 2010.

\bibitem{lawvere1969}
F.~W.~Lawvere.
\newblock Diagonal arguments and cartesian closed categories.
\newblock In {\em Category Theory, Homology Theory and their Applications II}, pages 134--145.
\newblock Springer, 1969.

\bibitem{connes2024}
A.~Connes and C.~Consani.
\newblock The Riemann-Roch strategy: A spectral approach to the Riemann Hypothesis.
\newblock arXiv:2511.22755, 2024.

\bibitem{kemp2010}
A.~H.~Kemp and D.~S.~Quintana.
\newblock The relationship between mental and physical health: Insights from the study of heart rate variability.
\newblock {\em International Journal of Psychophysiology}, 89(3):288--296, 2013.

\bibitem{potters2020}
M.~Potters and J.-P.~Bouchaud.
\newblock {\em A First Course in Random Matrix Theory}.
\newblock Cambridge University Press, 2020.

\bibitem{summa}
J.~N.~K.
\newblock Summa of the Gap: The Complete Architecture of the Vault.
\newblock Unpublished manuscript, 2026.

\end{thebibliography}

% ============================================================
% APPENDIX
% ============================================================

\appendix
\section{Calibration Data}

\begin{table}[h]
\centering
\caption{Extended spectral measurements across nine architectures.}
\begin{tabular}{lccccc}
\toprule
Architecture & Depth & Width & $d$ & $C$ (mean) & $\alpha$ \\
\midrule
MLP-3-128 & 3 & 128 & 0.0 & 0.823 & 7.71 \\
MLP-5-256 & 5 & 256 & 0.0 & 0.810 & 7.68 \\
MLP-7-512 & 7 & 512 & 0.0 & 0.798 & 7.65 \\
Peak-3-256 & 3 & 256 & 0.25 & 0.772 & 7.63 \\
Peak-5-256 & 5 & 256 & 0.25 & 0.762 & 7.60 \\
Thresh-5-256 & 5 & 256 & 0.50 & 0.709 & 7.48 \\
Res-5-256 & 5 & 256 & 0.75 & 0.634 & 7.39 \\
FullRes-3-256 & 3 & 256 & 1.0 & 0.571 & 7.31 \\
FullRes-5-256 & 5 & 256 & 1.0 & 0.561 & 7.28 \\
\bottomrule
\end{tabular}
\end{table}

\section{Falsification Conditions}

For each claim in this paper, we specify the measurement that would falsify it:

\begin{enumerate}
    \item \textbf{Universal tail:} Measure $\alpha$ on an architecture where $\alpha \notin [6.5, 8.5]$. The invariance claim requires revision.
    \item \textbf{Tail-bulk decoupling:} Find an architecture where $\alpha$ co-varies with $C$. The decoupling claim is false.
    \item \textbf{Transformer flatness (prediction):} Measure $C > 0.10$ or $\sigma(\log\det) > 0.01$ across transformer layers. The fixed-point interpretation requires revision.
    \item \textbf{MMI violation:} Find a triple of activation subspaces where $I_3 \leq 0$. The claim that logdet is not holographic survives (failure of one counterexample does not establish holographicity).
\end{enumerate}

\end{document}
